\subsection{Two lecture blocks and two levels of depth}

\begin{corebox}[First 45 minutes: convergence of sequences]
The essential route is the definition of a sequence limit, elementary examples, and the role of monotonicity and boundedness. The Cauchy criterion and the monotone convergence theorem deepen this picture by explaining how convergence can be recognised without already knowing the limit.
\end{corebox}

\begin{corebox}[Second 45 minutes: from sequences to function limits]
Heine's theorem transfers convergence to functions. After that bridge, one-sided limits, infinite behaviour, continuity, and basic discontinuities can be read as questions about all admissible ways of approaching a point.
\end{corebox}

Formal criteria and existence theorems are marked as deeper structural results. Computational transformations are applications of the theory, not substitutes for the definition. A student should first understand what convergence means, then learn how a particular expression reveals it.
